\section{问题分析}

四个问题按“数据刻画—规律拟合—资源配置—演进预测”的顺序衔接：问题一产出域级质量评分
与 1M 规模下的配比响应面；问题二以经典标度律为骨架，判断质量能否以可检验的方式进入
标度律；问题三在问题二给出的损失函数上求解算力约束下的配置；问题四在真实排行榜数据上
检验规模与非规模因素的作用并预测前沿。贯穿全文的难点不在于拟合精度，而在于\textbf{识别}：
每一个数值结论都必须说明它由哪类数据支撑、在什么范围内成立。

\subsection{问题一的分析}

质量信号共 22 个指标，来自六个打分体系，量纲、方向与分布各不相同，且同一体系内的
指标往往高度相关。若直接等权平均，指标最多的体系（11 个规则类指标）会主导评分，
因此需要先把各指标映射到可比的秩尺度，再按体系平衡聚合。质量“冲突”不应只用负相关
来定义：两个指标在总体上正相关时，单条记录上仍可能给出相反评价，故冲突需要同时在指标
对与记录两个层面度量。

配比建模的难点有三。其一，17 维配比受单纯形约束，存在大量精确零分量，常用的对数比变换
在零点无定义；其二，训练集 A4/A5 与检验集 A6--A11 规模不同，其中 A6/A7 与训练集同为
1M 参数规模、A8/A9 是同一设计在 60M 规模的重复、A10/A11 则规模与设计同时改变，三者
检验的是不同的命题；其三，A13、A15 是估算参照而非观测，与之吻合不能证明外推正确。

\subsection{问题二的分析}

经典标度律 \eqref{eq:classic} 含不可约损失 $E$，对数线性最小二乘会把 $E$ 吸收进指数，
故需在原始参数上做稳健的非线性拟合。更关键的是数据性质：主拟合表 B1 被标注为观测数据，
但必须先检验其损失列是否带有真实训练噪声；若它与某一已发表公式逐点一致，拟合只能证明
估计器能恢复生成公式，不能证明真实模型的规模规律。外部证据只能来自独立的观测表（B4 剔除
与 B1 同族的行、B5 文献表）。

质量进入标度律的途径受两重限制：可观测的质量–损失联合数据只有半合成表 B6--B8；而问题一
给出的 $Q$ 是相对秩评分，二者之间没有第一方映射。因此广义标度律能否被采纳，需要事先设定
可证伪的采纳条件，而不是以拟合优度决定。

\subsection{问题三的分析}
% STAGE A: deliberately empty. The Q3 analysis is written only after T-010 is
% accepted and merged; nothing here may anticipate its method or results.

\subsection{问题四的分析}

排行榜数据是真实观测，但时间、规模与模型类型相互纠缠：同一时期既有预训练基础模型，也有
大量微调与合并模型；多数模型只有排行榜提交日期而无发布日期；训练算力元数据只覆盖极少数
模型。因此需要先冻结分析总体与时间口径，再在总体上分离“与参数规模相关”和“与规模无关”
的变化，并把以参数量为尺度的分解与以训练算力为尺度的分解区分开来：参数量增长不等于
训练算力增长。损失与得分之间的映射只能在可比的锚点上建立，映射误差必须进入依赖它的预测。
